\documentclass[aspectratio=169,11pt]{beamer}

% Packages
\usepackage[utf8]{inputenc}
\usepackage[T1]{fontenc}
\usepackage{lmodern}
\usepackage{microtype}
\usepackage{amsmath,amssymb}
\usepackage{booktabs}
\usepackage{graphicx}
\usepackage{tikz}
\usepackage{pgfplots}
\usepackage{xcolor}
\usepackage{ragged2e}
\usepackage{colortbl}
\usepackage{array}
\usepackage{hyperref}
\usepackage{adjustbox}

\usetikzlibrary{shapes.geometric, arrows.meta, positioning, calc, backgrounds, decorations.pathreplacing, shadows, fadings}
\pgfplotsset{compat=1.18}

% --- Color Palette: Warm Professional ---
\definecolor{DeepNavy}{HTML}{2E4057}
\definecolor{Teal}{HTML}{048A81}
\definecolor{WarmOrange}{HTML}{E85D04}
\definecolor{SoftPurple}{HTML}{9D4EDD}
\definecolor{WarmGray}{HTML}{6C757D}
\definecolor{LightGray}{HTML}{E9ECEF}
\definecolor{Cream}{HTML}{FBF8F1}
\definecolor{DeepRed}{HTML}{D62828}
\definecolor{Gold}{HTML}{D4A03A}
\definecolor{SoftWhite}{HTML}{FAFAFA}

% Beamer colors
\setbeamercolor{normal text}{fg=DeepNavy, bg=SoftWhite}
\setbeamercolor{structure}{fg=DeepNavy}
\setbeamercolor{alerted text}{fg=DeepRed}
\setbeamercolor{example text}{fg=Teal}
\setbeamercolor{frametitle}{fg=DeepNavy, bg=SoftWhite}
\setbeamercolor{title}{fg=DeepNavy}
\setbeamercolor{subtitle}{fg=WarmGray}
\setbeamercolor{block title}{bg=DeepNavy, fg=SoftWhite}
\setbeamercolor{block body}{bg=LightGray, fg=DeepNavy}
\setbeamercolor{itemize item}{fg=WarmOrange}
\setbeamercolor{itemize subitem}{fg=Gold}

% Fonts
\usefonttheme{professionalfonts}
\setbeamerfont{title}{size=\huge, series=\bfseries}
\setbeamerfont{frametitle}{size=\Large, series=\bfseries}

% Strip chrome
\setbeamertemplate{navigation symbols}{}
\setbeamertemplate{headline}{}

% Minimal footline
\setbeamertemplate{footline}{%
    \hfill
    \begin{beamercolorbox}[wd=3cm,ht=2.5ex,dp=1ex,right,rightskip=0.6cm]{page number}
        \usebeamercolor[fg]{WarmGray}\scriptsize\insertframenumber
    \end{beamercolorbox}
    \vspace{0.3cm}
}

% Bullets
\setbeamertemplate{itemize item}{\footnotesize\raisebox{0.3ex}{\tikz\fill[WarmOrange] (0,0) circle (0.45ex);}}
\setbeamertemplate{itemize subitem}{\footnotesize\raisebox{0.3ex}{\tikz\fill[Gold] (0,0) circle (0.35ex);}}

\begin{document}

\begin{frame}
    \title{Appendix}
    \subtitle{GIV in Practice: PCA and GMM Implementation}
    \author{Roy Kisluk}
    \date{}
    \titlepage
\end{frame}

% --- SLIDE 1: RECAP ---
\begin{frame}
    \frametitle{The Need for $\Gamma$: Leaky Instruments}
    
    \vspace{0.5cm}
    {\Large
        If large firms respond differently to macro trends, the simple GIV fails.
    }
    
    \vspace{1cm}
    \begin{itemize}
        \item \textbf{The Problem:} If big firms are more sensitive to macro shocks ($\eta_t$), then $y_{St} - y_{Et}$ is correlated with the error term.
        \item \textbf{The Solution:} We need weights $\Gamma$ such that $\sum \Gamma_i = 0$ AND $\sum \Gamma_i \lambda_i = 0$.
        \item \textbf{The Formula:} $\Gamma = S - \lambda (\lambda' \lambda)^{-1} \lambda' S$ (the residual of $S$ on $\lambda$).
    \end{itemize}
\end{frame}

% --- SLIDE 2: PCA METHOD ---
\begin{frame}
    \frametitle{Implementation 1: PCA (The Two-Step Approach)}
    
    {\large
        Simplest when you have a long panel and common factors dominate.
    }
    
    \vspace{0.5cm}
    \textbf{Algorithm:}
    \begin{enumerate}
        \item \textbf{De-mean:} Center the outcome matrix $Y$ (firm-by-time).
        \item \textbf{Factorize:} Run PCA on $Y$ and extract the first few loadings $\hat{\lambda}_i$.
        \item \textbf{Project:} Calculate $\Gamma_i$ as the residuals of regressing size $S_i$ on $\hat{\lambda}_i$.
        \item \textbf{Instrument:} Construct $z_t = \sum \Gamma_i y_{it}$ and use in standard 2SLS.
    \end{enumerate}
    
    \vspace{0.5cm}
    \textit{Pro: Computationally cheap. Con: Ignores the endogenous variable's impact on factors.}
\end{frame}

% --- SLIDE 3: GMM METHOD ---
\begin{frame}
    \frametitle{Implementation 2: Iterative GMM (The Robust Approach)}
    
    {\large
        Simultaneously identifies the elasticity ($\psi$) and the loadings ($\lambda_i$).
    }
    
    \vspace{0.5cm}
    \textbf{Iterative Algorithm:}
    \begin{itemize}
        \item \textbf{Start:} Assume $\psi^{(0)} = 0$ and $\lambda_i^{(0)} = 1$.
        \item \textbf{Loop:}
        \begin{enumerate}
            \item \textbf{Residuals:} Calculate $e_{it} = y_{it} - \psi^{(k)} x_{it}$.
            \item \textbf{Loadings:} Extract $\lambda_i^{(k+1)}$ from $e_{it}$ (e.g., via PCA on residuals).
            \item \textbf{GIV:} Update $z_t = \sum \Gamma(\lambda^{(k+1)}) y_{it}$.
            \item \textbf{Estimate:} Update $\psi^{(k+1)}$ via IV regression of $y$ on $x$ using $z_t$.
        \end{enumerate}
        \item \textbf{Converge:} Stop when $\psi$ and $\lambda$ stabilize.
    \end{itemize}
\end{frame}

% --- SLIDE 4: SUMMARY ---
\begin{frame}
    \frametitle{Summary: When to use which?}
    
    \vspace{0.5cm}
    \begin{center}
    \renewcommand{\arraystretch}{1.5}
    \begin{tabular}{l l l}
        \toprule
        \textbf{Method} & \textbf{Complexity} & \textbf{Ideal Use Case} \\
        \midrule
        Simple GIV & Low & Homogeneous firms ($\lambda_i \approx \bar{\lambda}$) \\
        \rowcolor{Teal!10} PCA GIV & Medium & High Signal-to-Noise, $T$ is large \\
        \rowcolor{WarmOrange!10} GMM GIV & High & Endogeneity in factor structure \\
        \bottomrule
    \end{tabular}
    \end{center}
    
    \vspace{0.5cm}
    \textbf{Bottom Line:} In granular economies, \textbf{always} check if results are robust to using $\Gamma$ (the "Purified" GIV weights).
\end{frame}

\end{document}
